\section{Schaltgeräte und Schaltanlagen}
\subsection{Schaltgeräte}

    \subsubsection{Lichtbogen}
        \begin{itemize}
            \item entsteht wenn anliegende Spannung größer als die Zündspannung \(U_Z\) ist, 
                bleibt allerdings bestehen solange die Spannung größer der Brennspannung \(U_B\) ist. (\(U_B < U_Z\))
            \item Gefahr des Kontaktabbrands und der Verschmelzung der Kontakte.
            \item Da Lichtbogen durch Ionisation entsteht, ist die Umgebung ausschlaggebend für die Zündspannung \(U_Z\). 
            \item Lichtbogen wird ``kalt'' durch:
                \begin{itemize}
                    \item DC: Kühlung
                    \item AC: Nulldurchgang des Stroms
                \end{itemize}
        \end{itemize}
        
        \textbf{Ayrton-Gleichung}
        \begin{align*}
            U_B = a + b \cdot l + \frac{c + d \cdot l}{I}
        \end{align*}
        \(a, b, c, d\) : empirische Konstanten. Für Kupfer:
        \begin{table}[H]
            \begin{tabular}{l l}
                \toprule
                \(a\) & \SI{15}{\volt}  (Anode \SI{5}{\volt}, Kathode \SI{10}{\volt}) \\
                \(b\) & \SI{10}{\volt\per\centi\meter} \\
                \(c\) & \SI{10}{\volt\ampere} \\
                \(d\) & \SI{50}{\volt\ampere\per\centi\meter} \\
                \bottomrule
            \end{tabular}
        \end{table}
        {\small Diagramm: Kathode auf \SI{0}{\volt}}

